\documentclass[a4paper]{article}

% Hỗ trợ tiếng Việt Unicode
\usepackage[utf8]{vntex}

% KÍCH HOẠT POSTCORRECT: Thêm tùy chọn 'postcorrect' vào gói automultiplechoice
\usepackage[box,completemulti,postcorrect]{automultiplechoice}    

\begin{document}

% Số lượng bản đề thi: Tạo 3 bản
% Bản 1 (Sheet 1): Dành cho Giáo viên tô đáp án đúng làm bài mẫu (Answer key)
% Bản 2, Bản 3 (Sheet 2, 3): Dành cho Học sinh làm bài
%
% QUAN TRỌNG: Với postcorrect, KHÔNG xáo trộn câu hỏi/đáp án (không dùng \shufflegroup)
% để cấu trúc các bản đề hoàn toàn trùng khớp với bài chấm mẫu.
\onecopy{3}{    

\noindent{\bf BÀI KIỂM TRA ĐỊNH KỲ \hfill POSTCORRECT}

\vspace*{.5cm}
\begin{minipage}{.4\linewidth}
\centering\large\bf Đề thi trắc nghiệm\\ Môn: Kiến thức chung\end{minipage}
\namefield{\fbox{    
                \begin{minipage}{.5\linewidth}
                  Họ và tên:

                  \vspace*{.5cm}\namefielddots   
                  \vspace*{1mm}
                \end{minipage}
         }}

\vspace{1ex}

% --- CHÚ Ý ĐẶC BIỆT VỀ ĐÁP ÁN: ---
% Trong chế độ postcorrect, file LaTeX KHÔNG CẦN và KHÔNG DÙNG đáp án đúng.
% Tất cả các lựa chọn bên dưới ĐỀU ĐƯỢC ĐẶT LÀ \wrongchoice!
% Đáp án đúng thực sự sẽ do Giáo viên cầm bút tự tô lên Tờ số 1 (Sheet 1).

% --- CÂU HỎI 1 ---
\begin{question}{q1}    
  Thủ đô của Việt Nam là thành phố nào?
  \begin{choices}
    \wrongchoice{Hà Nội}
    \wrongchoice{TP. Hồ Chí Minh}
    \wrongchoice{Đà Nẵng}
    \wrongchoice{Huế}
  \end{choices}
\end{question}

% --- CÂU HỎI 2 ---
\begin{question}{q2}    
  Số nguyên tố chẵn duy nhất là số nào?
  \begin{choices}
    \wrongchoice{0}
    \wrongchoice{1}
    \wrongchoice{2}
    \wrongchoice{4}
  \end{choices}
\end{question}

% --- CÂU HỎI 3: Câu hỏi nhiều lựa chọn ---
% Khi chấm điểm, nếu Giáo viên tô 2 đáp án (Trung Quốc và Lào) trên tờ bài mẫu,
% AMC sẽ tự động coi 2 đáp án đó là đáp án đúng!
\begin{questionmult}{q3}    
  Những quốc gia nào giáp biên giới trên đất liền với Việt Nam?
  \begin{choices}
    \wrongchoice{Trung Quốc}
    \wrongchoice{Lào}
    \wrongchoice{Campuchia}
    \wrongchoice{Thái Lan}
  \end{choices}
\end{questionmult}

}

\end{document}
